\documentclass[../../../include/open-logic-section]{subfiles}

\begin{document}

\olfileid{sol}{syn}{exp}
\olsection{表达力}

\begin{explain}
对二阶变元的量化使语言的表达力相比一阶逻辑有了极大的提升。
二阶存在量化允许我们断言存在具有某种性质的函数或关系。
在一阶逻辑中，做到这一点的唯一方法就是为此指定一个非逻辑符号（即函数符号或谓词符号）。
二阶全称量化允许我们断言，论域的所有子集、论域上的所有关系，或从论域到论域的所有函数都具有某个性质。
在一阶逻辑中，我们只能说分配给语言中某个非逻辑符号的子集、关系或函数具有该性质。
而且，当我们断言存在具有某种性质的子集、关系或函数，或者断言所有子集、关系、函数都具有该性质时，我们在刻画该性质时同样可以使用二阶量化。
这使得我们能够定义一阶逻辑无法定义的关系，并表达一阶逻辑无法表达的论域性质。
\end{explain}

% Part: second-order-logic
% Chapter: syntax-and-semantics
% Section: expressive-power

\begin{defn}
如果 $\Struct{M}$ 是语言~$\Lang{L}$ 的一个结构，
关系~$R \subseteq \Domain{M}^2$ 在~$\Lang{L}$ 中是\emph{可定义的}，
当且仅当存在某个公式~$!A_R(x, y)$，其中仅变元 $x$ 和~$y$
自由，使得 $R(a, b)$ 成立（即 $\tuple{a,b} \in R$）
当且仅当对满足 $s(x) = a$ 且 $s(y)=b$ 的 $s$，有 $\Sat{M}{!A_R(x, y)}[s]$。
\end{defn}

\begin{ex}
在一阶逻辑中，我们可以用公式 $\eq[x][y]$ 来定义等词关系~$\Id{\Domain{M}}$
（即 $\Setabs{\tuple{a,a}}{a \in
  \Domain{M}}$）。
在二阶逻辑中，我们\emph{无需~$\eq$} 就能定义这个关系。
因为如果 $a$ 和 $b$ 是~$\Domain{M}$ 中的同一个元素，
那么它们就属于~$\Domain{M}$ 的相同子集（因为集合由其元素决定）。
反之，如果 $a$ 和 $b$ 不同，那么它们就不属于完全相同的子集：
例如，若 $a \neq b$，则 $a \in \{a\}$ 而 $b \notin \{a\}$。
因此，“属于~$\Domain{M}$ 的相同子集”是这样一个关系：它在 $a$ 和 $b$ 之间成立，当且仅当 $a = b$。
这是一个可以在二阶逻辑中表达的关系，因为我们可以对~$\Domain{M}$ 的所有子集进行量化。
因此，下列公式定义了 $\Id{\Domain{M}}$：
\[
\lforall[X][(X(x) \liff X(y))]
\]
\end{ex}

\begin{prob}
证明 $\lforall[X][(X(x) \lif X(y))]$（注意：是 $\lif$ 而非~$\liff$！）定义了 $\Id{\Domain{M}}$。
\end{prob}

\begin{ex}
若 $R$ 是一个二元谓词符号，则 $\Assign{R}{M}$ 是~$\Domain{M}$ 上的一个二元关系。
可能有点令人困惑的是，我们将用 $R$ 同时表示谓词符号~$R$ 和关系~$\Assign{R}{M}$ 本身。
$R$ 的\emph{传递闭包}~$R^*$ 是这样一个关系：$R^*(a,b)$ 成立当且仅当存在 $c_1$, \dots, $c_k$ 使得
$R(a,c_1)$, $R(c_1, c_2)$, \dots, $R(c_k,b)$ 成立。
这包括 $k = 0$ 的情形，即如果 $R(a,b)$ 成立，则 $R^*(a,b)$ 也成立。
这意味着 $R \subseteq R^*$。事实上，$R^*$ 是包含~$R$ 且传递的最小关系。
我们可以在二阶逻辑中说 $X$ 是一个包含~$R$ 的传递关系：
\begin{multline*}
  !B_R(X) \ident \lforall[x][\lforall[y][(R(x,y) \lif X(x, y))]] \land {}\\
\lforall[x][\lforall[y][\lforall[z][((X(x,y) \land X(y,z)) \lif X(x,
      z))]]].
\end{multline*}
第一个合取支说 $R \subseteq X$，第二个则说 $X$ 是传递的。

说 $X$ 是这样的最小关系，就是说 $X$ 本身包含于每一个包含 $R$ 且传递的关系中。
因此，我们可以用以下公式定义 $R$ 的传递闭包：
\[
R^*(X) \ident !B_R(X) \land \lforall[Y][(!B_R(Y) \lif
  \lforall[x][\lforall[y][(X(x, y) \lif Y(x,y))]])].
\]
我们有 $\Sat{M}{R^*(X)}[s]$ 当且仅当 $s(X) = R^*$。
$R$ 的传递闭包无法在一阶逻辑中表达。
\end{ex}

\end{document}